\documentclass{oxmathproblems}

\printanswers

\course{Understanding Analysis, 2nd Edition}
\sheettitle{Chapter 5 - Section 4.6 - Sets of Discontinuity} 

\begin{document}
\begin{questions}

%------------------------- Problem 1 -------------------------
\miquestion Define $D_{f} \subset \textbf{R}$ to be the set of points where the function $f$ is discontinuous. Using modifications of Dirichlet / Thomae functions,
construct a function $f: \textbf{R} \to \textbf{R}$ so that
\begin{enumerate}[label=(\alph*)]
  \item $D_{f} = \textbf{Z}^c$.
  \item $D_{f} = \{x: 0 < x \leq 1\}$.
\end{enumerate}

\begin{solution}
  \begin{enumerate}[label=(\alph*)]
    \item Let $D(x)$ be the distance from $x$ to the integer closest to it. That is, $D(x) = \min(|x - \lfloor x \rfloor|, |x - \lceil x \rceil|)$.
      \[f(x) = 
      \begin{cases}
         0 & \text{if } x \text{ is an integer or irrational} \\
         D(x) & \text{if } x \text{ is rational}
      \end{cases}
     \]

     It's continuous at integer $x$. If $(x_{n}) \to x$ then $f(x_{n}) \to 0$ for both rationals and irrationals in $(x_{n})$.

     It's discontinuous at rational $x$. Take a sequence of irrationals $(x_{n}) \to x$. We have $f(x_{n}) \to 0$ but $f(x) > 0$.

     It's discontinuous at irrational $x$. Take a sequence of rationals $(x_{n}) \to x$. We have $f(x_{n}) \to D(x) > 0$ but $f(x) = 0$.
     
    \item Rename the function in part(a) to $g$. We'll reuse it here. Let
    \[f(x) = 
      \begin{cases}
         0 & \text{if } x \leq 0 \\
         g(x) & \text{if } 0 < x < 1 \\
         1 & \text{if } x \geq 1
      \end{cases}
    \]

    It's continuous on $x < 0$ because it's constant there.\\
    It's continuous at 0 because by part (a) $x \to 0 \implies g(x) \to 0$ and $g(0)=0$.\\
    It's discontinuous at $0 < x < 1$ by part(a).\\
    It's discontinuous at 1 because by part(a), $x^{-} \to 1 \implies g(x) \to 0$ but $g(1)=1$.\\
    It's continuous on $x > 1$ because it's constant there.
  \end{enumerate}
\end{solution}

%------------------------- Problem 2 -------------------------
\miquestion Given a countable set $A = \{a_{1}, a_{2}, a_{3}, ...\}$, define $f(a_{n}) = 1/n$ and $f(x)=0$ for all $x \notin A$. Find $D_{f}$.
\begin{solution}
  It's not too difficult to visualize that $f$ must be discontinuous at $c \in A$. Any $\delta$-neighborhood  $V_{\delta}(c)$ must contain an
  $x \notin A$ (if not then there exists an interval that consists only of $x \in A$, making $A$ uncountable). Then any $V_\delta(c)$ contains an
  $x$ such that $f(x)=0$ and since $f(c) > 0$, it's discontinuous.

  What about $c \notin A$? If $A=\textbf{Z}$ then it's clearly continuous. If $A = \textbf{Q}$ it's less obvious, but turns out it's still continuous.

  Given $\epsilon > 0$, let $N > \frac{1}{\epsilon}$. Since $N$ is bounded, we can take $m = \min_{1 \leq i \leq N}|a_{i}-c|$.
  Then for $\delta < m$, $V_\delta(c)$ contains only those $x$ where $f(x)=0$ if $x \notin A$ or $f(x) < \frac{1}{N} < \epsilon$ if $x \in A$.
  So $f$ is continuous at $c \notin A$.

  In conclusion, $D_{f}=A$.
\end{solution}

%------------------------- Problem 3 -------------------------
\miquestion Define the left-hand limit $\lim_{x \to c^{-}}f(x) = L$.
\begin{solution}
  For all $\epsilon > 0$, there exists a $\delta > 0$ such that $|f(x)-L| < \epsilon$ whenever $0< c-x < \delta$.
  Equivalently, in terms of sequences, $\lim_{x \to c^{-}}f(x) = L$ if $\lim f({x_{n}})=L$ for all sequences $(x_{n})$ satisfying $x_{n} < c$ and $\lim(x_{n})=c$.
\end{solution}

%------------------------- Problem 4 -------------------------
\miquestion Given $f: A \to \textbf{R}$ and a limit point $c$ of $A$, prove that $\lim_{x \to c}f(x) = L$ if and only if  $\lim_{x \to c^{-}}f(x) = L$ and  $\lim_{x \to c^{+}}f(x) = L$.
\begin{solution}
  Necessity: $\lim_{x \to c}f(x) = L$ means $(x_{n}) \to c \implies f(x_{n}) \to L$. This works regardless of whether $x_{n} > c$ or $x_{n} < c$.
  
  Sufficiency: Given $(x_{n}) \to c$, split it into two subsequences $(a_{n})$ for $x_{n} > c$ and $(b_{n})$ for $x_{n} < c$. Then $f(a_{n}) \to L$ and $f(b_{n}) \to L$
  so $f(x_{n}) \to L$ as well.
\end{solution}

%------------------------- Problem 5 -------------------------
\miquestion Prove that the only type of discontinuity a monotone function can have is a jump discontinuity.
\begin{solution}
  We first prove that a monotone function can have jump discontinuity. Pretty simple: $f(x) = 0$ for $x \leq 0$ and $f(x) = 1$ otherwise.

  Next we prove impossibility of removable discontinuity.  We'll prove it by contradiction. WLOG $f$ is increasing.
  
  Suppose $\lim_{x \to c}f(x)=L$ but $f(c) \neq L$. If $L > f(c)$, let $d=L-f(c)$. Then there exists $\delta$ such that
  $|x-c| < \delta \implies |f(x)-L| < \frac{d}{2}$. Choose $a < c$ with $|a-c| < \delta$, then $a < c$ but $f(a) > f(c)$, contradicting
  that $f$ is increasing.

  If $L < f(c)$, choose $b > c$ with $|b-c| < \delta$, then $b > c$ but $f(c) > f(b)$, again contradicting that $f$ is increasing.

  Finally, we prove impossibility of essential discontinuity. At first I thought the book is wrong, with $f(x)=1/x$ on $(0, 1]$ having
  essential discontinuity at $x=0$. But in this case $f(0)$ is undefined, hence the notion of continuity / discontinuity at that point
  is undefined as well.
  
  \textbf{----- Incorrect -----}\\
  Consider a limit point $c$ and a sequence $(a_{n}) \to c$ with $a_{n} < c$. Since $f(c)$ is defined and $f$ is monotone, $f(a_{n})$ is monotone and bounded.
  By the Monotone Convergence Theorem, $f(a_{n})$ converges. Similarly for a sequence $(b_{n})$ with $b_{n} > c$. Note that it's possible for $(a_{n})$ or $(b_{n})$
  to be empty, in which case it's vacuously true that they converge.

  \textbf{----- Correct -----}\\
  Two mistakes. (1) Just because $(a_{n}) \to c$ does not guarantee that $(a_{n})$ is monotone. (2) We need a single limit that works for all sequences to the left of $c$,
  not only a particular sequence $(a_{n})$.

  Instead we'll use supremum. Wlog $f$ is increasing. We'll prove the left limit exists.

  The set $\{f(x): x < c\}$ is bounded above by $f(c)$, so it has a supremum $L$. By definition of supremum, given an $\epsilon > 0$, there
  exists a point $a < c$ such that $L-f(a) < \epsilon$. Let $\delta =c-a$. Then by monotonicity, $0 < c-x < \delta \implies L-f(x) < \epsilon$. Hence the left limit is $L$.

  An analogous argument with infimum proves the right limit also exists.

  Since both the left and right limits exist, either they're equal hence the usual limit exists, or they're unequal in which case it's a jump discontinuity.
  Either case, an essential discontinuity is impossible.
\end{solution}

\end{questions}
\end{document}
